% -*- TeX-master: "../main.tex" -*-

\chapter{Conclusiones y trabajo futuro}

Este capítulo sintetiza los principales resultados obtenidos en el desarrollo del intérprete de GPLC y proyecta líneas de trabajo futuro relevantes. Por una parte, se presentan las conclusiones del trabajo, destacando el alcance del prototipo como implementación ejecutable de la semántica del lenguaje, su valor como medio de experimentación y la decisión de diseño de trasladar la formulación teórica fielmente a una implementación concreta. Por otra, se exponen líneas de trabajo futuro motivadas por las limitaciones identificadas durante la validación, incluyendo aspectos relacionados con la semántica formal, la representación del resultado, la resolución de fórmulas, el tipado de operaciones, la equivalencia semántica y la optimización de operaciones.



\section{Conclusiones}

\subsection{Implementación de la semántica de GPLC}



El principal resultado de este trabajo es la construcción de un intérprete que materializa, de manera ejecutable, la semántica estática y dinámica de GPLC presentada en la formulación de \cite{gplc} y la confirmación experimental de que dicha formulación funciona cuando es llevada a la práctica. Esto se traduce en dos hechos fundamentales: (1) el sistema de tipos rechaza estáticamente programas inconsistentes de acuerdo con el tipado gradual, que incorpora, además del tipo desconocido, probabilidades desconocidas, y (2) la semántica operacional reduce los programas a un resultado concreto y previsible según las reglas de la semántica operacional, y es capaz de rechazar dinámicamente programas con errores de tipo.

\subsection{Plataforma de experimentación}
El intérprete se constituye como un medio efectivo de experimentación
para la investigación relacionada con el lenguaje, lo cual cumple el
objetivo fundamental de este trabajo. Mediante el soporte que entrega
a la totalidad de sus constructos, es capaz de cumplir un rol
experimental relevante en la ejecución de programas de interés, en el
análisis del comportamiento del sistema de tipos, en la comprobación de
resultados esperados y en la verificación del término exitoso de programas o de errores controlados.


\subsection{Arquitectura alineada con la semántica formal}
El diseño de la arquitectura del intérprete tiene una correspondencia muy estrecha
con la estructura presentada en la exposición del lenguaje: el pipeline principal de
tareas refleja de manera transparente y directa las fases de tipado, elaboración y evaluación. Este hecho facilita la trazabilidad de la teoría en la implementación, permitiendo razonar con mayor facilidad acerca del diseño y funcionamiento interno del intérprete. Entre las ventajas que esto otorga, se encuentran una mayor facilidad de corrección y extensibilidad para otros desarrolladores que estén familiarizados con GPLC.

\subsection{Revisión de la teoría}
La implementación entregó evidencia concreta de casos en que la
ejecución produce resultados no esperados, o no termina. Este tipo de
fenómenos, además de motivar revisiones focalizadas de la
implementación, también fundamenta un examen crítico de reglas
específicas del formalismo. Durante el proceso de desarrollo, la
identificación de inconsistencias en el funcionamiento del intérprete
ya ha impulsado el rediseño de reglas de tipado en el cálculo formal
de GPLC, lo cual demuestra el potencial de la experimentación como
mecanismo de mejora de la teoría.


\subsection{Ineficiencias y cobertura de programas}
Finalmente, en el diseño del intérprete, el foco no se ha colocado en la
optimización de su funcionamiento, puesto que el objetivo esencial es
proveer un medio de experimentación para un modelo teórico. Sin
embargo, la ineficiencia asociada al manejo de operaciones con
distribuciones simbólicas limita la magnitud de los programas que el
intérprete puede ejecutar, y por lo tanto, restringe su alcance como
herramienta de investigación. Por lo tanto, hacen falta ajustes en el
diseño de algoritmos con recursión profunda y acumulación de fórmulas,
que actualmente están implementados como traducciones literales de la
formulación teórica.

\section{Trabajo futuro}

\subsection{Correcciones en la semántica formal}

 Como trabajo futuro, todo programa que demuestre un comportamiento
erróneo o no esperado justifica no solo un examen y pruebas acuciosas
de la implementación del código fuente del intérprete, sino también
una revisión de la semántica formal del lenguaje en que se basa
esta implementación. En particular, los
problemas relacionados con la regla de reducción de la ascripción
\lstinline|Dascr| observados en la ejecución fundamentan una revisión
de dicha regla.

\subsection{Representación del resultado} En el prototipo, el resultado se presenta como una distribución con pesos simbólicos, pero sin exponer la fórmula que restringe esos pesos. Esto puede volver la salida poco informativa cuando se requieren probabilidades numéricas o una explicación de las restricciones.

Un trabajo futuro natural es mejorar la forma en que el intérprete
comunica el resultado al usuario. En primer lugar, sería valioso
permitir que la ejecución muestre la fórmula que restringe las
probabilidades, de manera parcial y estratégica, y en un formato que
favorezca su legibilidad, pues la fórmula explica qué asignaciones
reales son compatibles con la distribución retornada.  En segundo
lugar, sería deseable habilitar un modo de instanciación de
probabilidades simbólicas, ya sea extrayendo un modelo desde el solucionador
cuando exista, o permitiendo que el usuario entregue valores para las
probabilidades desconocidas \lstinline|?|.



\subsection{Resolución de fórmulas} La satisfacibilidad de fórmulas se consulta de forma transversal y repetida durante la evaluación. Sin embargo, cuando cada invocación al solucionador se resuelve desde cero, el costo se acumula especialmente en operaciones donde las restricciones se forman por agregación de subfórmulas.

Un trabajo futuro relevante es explorar un enfoque incremental que
permita reutilizar información entre chequeos sucesivos, reduciendo el
costo de resolver sistemas que se forman, en parte, por acumulación de
restricciones ya chequeadas. Sin embargo, esta mejora debe diseñarse
con cuidado, dado que, incluso si una subfórmula es soluble, su
solución no necesariamente es compatible con una solución del sistema
completo cuando se agregan nuevas restricciones. Por ello, una
dirección razonable es reutilizar contexto y restricciones de manera
segura, sin asumir extensibilidad automática de
subfórmulas. Una ventaja en este sentido es que el solucionador Z3
  cuenta con soporte nativo para resolución incremental, manteniendo un
  contexto de restricciones que puede extenderse y revertirse. Sin
  embargo, este soporte no se traduce automáticamente en un enfoque
  incremental en el intérprete, pues requiere decidir qué contexto
  preservar, qué restricciones mantener y cómo reutilizarlas,
  sin asumir extensibilidad de soluciones parciales al agregar nuevas
  restricciones. En consecuencia, cualquier diseño de algoritmo o heurísticas
  con respecto a este problema puede apoyarse en las funcionalidades
  incrementales de Z3, pero integrándolas de manera segura.

\subsection{Tipado de operaciones binarias} El diseño
  actual del sistema de tipos interpreta las operaciones binarias como
  expresiones con tipo de distribución de Dirac, lo que entra en tensión con
  constructos que esperan valores simples. En consecuencia,
  subexpresiones aritméticas y booleanas pueden quedar conceptualmente
  ambiguas entre \emph{valor} y \emph{distribución}.

Esta limitación sugiere la necesidad de un rediseño que compatibilice
la interpretación de expresiones como distribuciones con los contextos
del lenguaje que esperan valores simples. Una línea de trabajo futuro
consiste en introducir un mecanismo explícito que permita tratar, por
ejemplo, \lstinline|{Real:1}| como \lstinline|Real| cuando el contexto
lo requiera, o bien adoptar un sistema de tipado más sensible al tipo
esperado por el contexto. En particular, uno de los desafíos está dado
por el hecho de que la elaboración transforma las operaciones binarias
en $\LetKW$ (el cual está típicamente asociado a la generación de
distribuciones), como se muestra a continuación:
\begin{equation*}
  v1 + v2  \mathrel{\leadsto} \LetTerm{x}{\eascr{\epsilon_1}{v_1}{\Real}}{\LetTerm{y}{\eascr{\epsilon_2}{v_2}{\Real}}{x+y}}
\end{equation*}
Sería importante, por lo tanto, distinguir
en el lenguaje destino entre los $\LetKW$ insertados
por la elaboración para chequear los valores
sumados, y aquellos que provienen de un $\LetKW$ en el
lenguaje fuente.

\subsection{Equivalencia semántica de expresiones}
 Actualmente no existe un mecanismo interno para comparar
  expresiones o distribuciones resultantes. Esto dificulta automatizar
  pruebas que contrasten formalmente un resultado esperado con el
  observado, y limita la validación sistemática del comportamiento del
  lenguaje.

Una extensión importante del lenguaje sería
incorporar un mecanismo para comparar expresiones
del lenguaje fuente por equivalencia semántica, lo
que se traduce en comparar las distribuciones de
valores que dichas expresiones generan. Para ello
será necesario definir primero qué significa
igualdad entre valores (incluyendo el rol del tipo
ascrito) y decidir cómo tratar valores
funcionales, que se evalúan a clausuras. Con esa
base, podría definirse una noción de igualdad
entre distribuciones que considere tanto el
soporte como las restricciones sobre
probabilidades simbólicas. En presencia de
probabilidades desconocidas, también resulta natural estudiar
una igualdad parcial: distribuciones
que coinciden para algún conjunto no vacío de
instanciaciones de probabilidades simbólicas.




\subsection{Optimización de operaciones}  En el lenguaje
  destino, \lstinline|Let| y \lstinline|Dascr| aparecen
  recurrentemente y ejecutan operaciones costosas entre distribuciones
  de tipos en tiempo de ejecución. En programas no triviales, este
  patrón eleva sustancialmente el costo y restringe la escalabilidad
  del prototipo.  Por ello, un trabajo futuro relevante es estudiar
  estrategias para reducir el crecimiento del tamaño de distribuciones
  y fórmulas durante la evaluación. En términos generales, esto apunta
  a evitar combinaciones innecesarias en operaciones $m\times n$, y a
  simplificar distribuciones y restricciones antes de delegar al
  solucionador. El objetivo sería mejorar la escalabilidad del prototipo,
  permitiendo ejecutar programas más grandes y observar
  comportamientos que hoy quedan fuera por costo computacional.

